% !TEX root = ../../main.tex

\section{Cumplimiento de Objetivos}

La Tabla siguiente resume el cumplimiento de los objetivos definidos en la Introducción.

\begin{center}
\begin{tabular}{p{0.46\textwidth}p{0.40\textwidth}}
\hline
Objetivo & Resultado \\
\hline
Analizar \texttt{ApplicativeDo} en \textsf{GHC}. & Se identificó el flujo del Renamer para \texttt{HsDo}, el punto donde se procesa \texttt{ApplicativeDo} y la construcción de statements aplicativos. \\
Diseñar un modelo de dependencias locales. & Se definió un grafo RAW usando variables definidas, variables libres y lecturas locales por statement. \\
Implementar la extensión experimental en \textsf{GHC}. & Se integraron el reconocimiento de la marca explícita, la generación de permutaciones topológicas, la evaluación de candidatos y la selección automática o forzada dentro del flujo de \texttt{ApplicativeDo}. \\
Evaluar la preservación de resultados y la reducción potencial del costo estructural. & Se probaron los corpus sintéticos de \texttt{Maybe} y \texttt{Dist.T Rational} con logs reproducibles de costos, grafos y salidas; además, el control negativo con \texttt{IO} mostró el riesgo de una declaración de conmutatividad incorrecta. \\
\hline
\end{tabular}
\end{center}

En conjunto, los objetivos quedaron cubiertos por la modificación de \textsf{GHC}, los módulos agregados, los dos corpus sintéticos, el control negativo con \texttt{IO}, los scripts de validación y la documentación del repositorio.
